\chapter{System Services}
\label{ch:services}

Veilbox runs four background services: DHCP client (udhcpc), system logger
(syslogd), container runtime (containerd), and SSH server (Dropbear). This
chapter documents each service's configuration, operation, and integration.

\section{Service Architecture}

All services are started by \texttt{/etc/init.d/rcS} during the sysinit
phase of BusyBox init. They run in the background and are managed through
simple shell scripts --- there is no systemd, no supervisor, no service
manager.

\begin{figure}[H]
\centering
\begin{tikzpicture}[node distance=1cm]
\node[highlight] (init) {BusyBox Init (PID 1)};
\node[box, below left of=init, xshift=-2cm] (rcS) {/etc/init.d/rcS};
\node[box, below of=rcS, xshift=-1.5cm] (dhcp) {udhcpc};
\node[box, below of=rcS, xshift=0cm] (syslog) {syslogd};
\node[box, below of=rcS, xshift=1.5cm] (ctrd) {containerd};
\node[box, below of=rcS, xshift=3cm] (drop) {Dropbear};
\draw[arrow] (init) -- (rcS);
\draw[arrow] (rcS) -- (dhcp);
\draw[arrow] (rcS) -- (syslog);
\draw[arrow] (rcS) -- (ctrd);
\draw[arrow] (rcS) -- (drop);
\end{tikzpicture}
\caption{Service startup hierarchy. BusyBox init runs rcS, which forks
four background processes.}
\label{fig:services}
\end{figure}

\section{udhcpc: DHCP Client}

\subsection{Purpose}

The BusyBox DHCP client (\texttt{udhcpc}) negotiates an IPv4 address from
the network's DHCP server. In QEMU user-mode networking, the built-in SLiRP
DHCP server assigns 10.0.2.15/24 with gateway 10.0.2.2.

\subsection{Configuration}

The DHCP client requires a script that configures the interface when an
address is obtained. Veilbox provides:

\begin{lstlisting}[language=bash]
# /usr/share/udhcpc/default.script
#!/bin/sh
case "$1" in
    bound|renew)
        ifconfig $interface $ip netmask $subnet
        route add default gw $router
        echo "DNS: $dns" > /etc/resolv.conf
        ;;
    deconfig)
        ifconfig $interface 0.0.0.0
        ;;
esac
\end{lstlisting}

\subsection{Operation}

The script defines three actions:
\begin{itemize}[nosep]
  \item \texttt{bound}: Called when a lease is first acquired. Configures
        the IP address, netmask, default route, and DNS.
  \item \texttt{renew}: Called when a lease is renewed. Re-applies
        network settings.
  \item \texttt{deconfig}: Called when the interface is brought down.
        Removes the IP address.
\end{itemize}

\subsection{Troubleshooting}

Common DHCP issues:
\begin{itemize}[nosep]
  \item \texttt{udhcpc: ifconfig failed}: The interface name does not
        match. Verify with \texttt{ip link show}. Default is \texttt{eth0}.
  - \texttt{No DHCPOFFERS received}: The virtual network is not configured.
        Add \texttt{-netdev user,...} to QEMU command line.
  - \texttt{Lease failed, renewing in 60s}: The DHCP server is slow or
        unreachable. Verify QEMU's SLiRP is running.
\end{itemize}

\section{syslogd: System Logger}

\subsection{Purpose}

BusyBox syslogd captures kernel log messages (from /proc/kmsg) and process
messages (from /dev/log) and writes them to /var/log/messages.

\subsection{Configuration}

\begin{lstlisting}[language=bash]
syslogd -O /var/log/messages -s 64
\end{lstlisting}

Options:
\begin{itemize}[nosep]
  \item \texttt{-O}: Output file path.
  \item \texttt{-s}: Maximum message size (64 KB).
\end{itemize}

\subsection{Log Format}

Each log entry follows the standard syslog format:
\begin{lstlisting}
Jun  7 10:15:30 veilbox daemon.info containerd[42]: starting...
Jun  7 10:15:31 veilbox auth.info dropbear[58]: Running as server
\end{lstlisting}

\section{containerd: Container Runtime}

\subsection{Purpose}

Containerd is an industry-standard container runtime that manages the
complete container lifecycle: image pull, storage, execution, and
monitoring. It communicates via a gRPC API over a Unix domain socket.

\subsection{Configuration}

\begin{lstlisting}
# /etc/containerd/config.toml
version = 2
root = "/mnt/state/containerd"
state = "/run/containerd"

[grpc]
  address = "/run/containerd/containerd.sock"
  uid = 0
  gid = 0

[metrics]
  address = "127.0.0.1:1338"

[plugins]
  [plugins."io.containerd.grpc.v1.cri"]
    sandbox_image = "registry.k8s.io/pause:3.9"
\end{lstlisting}

\subsection{Startup}

\begin{lstlisting}[language=bash]
containerd -c /etc/containerd/config.toml &
\end{lstlisting}

\subsection{Interaction}

Use nerdctl (or the docker alias) to interact with containerd:
\begin{lstlisting}
nerdctl images                   # List pulled images
nerdctl pull alpine:latest       # Pull an image
nerdctl run -it alpine:latest sh # Run a container
nerdctl ps -a                    # List all containers
nerdctl logs <id>               # View container logs
\end{lstlisting}
The \texttt{docker} command is a symlink to \texttt{nerdctl}, so every
example above also works with \texttt{docker images}, \texttt{docker pull},
etc. No Docker daemon is needed inside the VM.

\section{Dropbear: SSH Server}

\subsection{Purpose}

Dropbear is a lightweight SSH server (approximately 140 KB) that provides
encrypted remote access to the Veilbox shell. It supports key-based and
password authentication.

\subsection{Configuration}

\begin{lstlisting}[language=bash]
dropbear -F -R -p 22
\end{lstlisting}

Options:
\begin{itemize}[nosep]
  \item \texttt{-F}: Run in foreground (for rcS).
  \item \texttt{-R}: Automatically create host keys if missing.
  \item \texttt{-p 22}: Listen on port 22.
  \item \texttt{-w}: Disallow root password login (not used in Veilbox).
\end{itemize}

\subsection{Host Keys}

Dropbear generates two host keys during build:
\begin{itemize}[nosep]
  \item \texttt{/etc/dropbear/dropbear\_ed25519\_host\_key}: ED25519 key
        (fast, modern, recommended).
  \item \texttt{/etc/dropbear/dropbear\_rsa\_host\_key}: RSA 2048-bit
        key (fallback for older clients).
\end{itemize}

\subsection{Authentication Methods}

\begin{itemize}[nosep]
  \item \textbf{Password}: User \texttt{root} with password \texttt{veiladmin}.
        The password is hashed with SHA-512 in /etc/shadow.
  \item \textbf{Public key}: ED25519 key in
        \texttt{/etc/dropbear/authorized\_keys}. The private key is
        provided in the output directory.
\end{itemize}

\subsection{Connecting}

\begin{lstlisting}[language=bash]
ssh -i output/ssh-test-key root@localhost -p 2222
# Or with password:
ssh root@localhost -p 2222
# Password: veiladmin
\end{lstlisting}

\section{Service Logging}

All service output can be monitored:
\begin{lstlisting}
cat /var/log/messages          # System log
nerdctl logs <container-id>   # Container logs
dropbear -F -E                 # Verbose SSH debug
\end{lstlisting}
